%==========================================================================
% CAPÍTULO 7 — CONCLUSÕES E TRABALHO FUTURO
%==========================================================================
\chapter{Conclusões e Trabalho Futuro}

    Este capítulo encerra o relatório com uma abordagem crítica ao trabalho realizado ao longo das quatro etapas do projeto e com a identificação de linhas de orientação para a sua evolução. O objetivo não é repetir o que já foi descrito nos capítulos anteriores, mas sim olhar para o conjunto, destacar o que merece ser preservado em projetos futuros desta natureza e reconhecer com franqueza as limitações da solução produzida.

%--------------------------------------------------------------------------
\section{Conclusões}
%--------------------------------------------------------------------------

    O projeto Belo Mercado Manager foi conduzido ao longo de quatro etapas alinhadas com um ciclo de vida clássico de Engenharia de \textit{Software} (Sommerville, 2016), enriquecido com a utilização sistemática de \textit{Large Language Models} (LLM) como agentes de apoio. A Etapa~1 consolidou a base de requisitos: 19 requisitos funcionais, 16 requisitos não funcionais, 29 \textit{user stories} e 19 casos de uso detalhados segundo o formato expandido de Cockburn, com alinhamento à norma IEEE~830 e ISO/IEC/IEEE~29148. A Etapa~2 traduziu essa base num desenho técnico: arquitetura em camadas com instância única centralizada (Blazor/.NET, SQL~Server, JWT, API REST), suportada por dez \textit{Architecture Decision Records}, diagramas de classes, sequência e componentes, modelo lógico de dados e seis \textit{wireframes}. A Etapa~3 implementou a solução em seis projetos .NET~8 com cerca de 22 entidades de domínio, 10 serviços de aplicação e respetivos repositórios, uma API REST documentada com \textit{Swagger}, e dois \textit{front-ends} Blazor Server. A Etapa~4 fechou o ciclo com uma \textit{suite} de 165 testes automáticos a passar de forma reprodutível, cobertura total dos requisitos funcionais em âmbito e uma avaliação ISO/IEC~25010 com nota média de 3,8 em 5.

    A primeira leitura crítica que importa fazer respeita à \textbf{adequação do método ao prazo disponível}. A organização do trabalho em quatro etapas claramente delimitadas, cada uma com artefactos próprios e marcos de validação, foi determinante para entregar um produto coerente em pouco mais de três meses. A rastreabilidade explícita entre requisitos, casos de uso, diagramas, código e testes (formalizada em matriz no Capítulo~6) tornou possível responder, em qualquer momento, à pergunta ``onde é que este requisito está implementado e validado?''. Esta disciplina documental compensou o pouco tempo disponível para cada etapa individualmente.

    Em segundo lugar, sobre a \textbf{utilização do LLM}, os resultados foram concretamente mensuráveis. Na geração de código, o LLM revelou-se especialmente produtivo a tratar das partes repetitivas (mapeamentos DTO/entidade, repositórios, controladores) e a aplicar correções cirúrgicas quando os \textit{bugs} eram descritos com precisão. Na geração de testes, a sua capacidade de propor \textit{edge cases} surpreendeu pela positiva: o caso do NIF \texttt{"000000000"}, matematicamente válido pelo algoritmo \textit{mod-11} mas contra-intuitivo, é o exemplo mais ilustrativo. A produtividade observada na fase de testes (cerca de 4$\times$ face à escrita manual estimada) sugere que esta é uma das áreas de maior potencial.

    Em terceiro lugar, importa reconhecer os \textbf{limites encontrados na utilização do LLM}. O modelo tende a inventar APIs plausíveis quando não tem os ficheiros relevantes carregados em contexto, o que originou erros recorrentes (assinaturas erradas de repositórios, nomes de \textit{enum} inexistentes) só apanhados em compilação. Tem também um padrão de \textit{otimismo de baseline} em avaliações qualitativas: a primeira proposta de notas ISO/IEC~25010 atribuía 5 em quase todas as características, claramente desalinhada com a maturidade real do produto. Finalmente, e mais importante, o LLM não consegue avaliar de forma autónoma a correção de uma regra de negócio: revê estilo e estrutura mas não questiona uma fórmula de cálculo errada. Este facto motivou diretamente a estratégia de testes adotada na Etapa~4, em que cada regra crítica é confirmada por uma asserção explícita.

    Por fim, sobre o \textbf{produto entregue}, a avaliação ISO/IEC~25010 (ISO/IEC, 2011) é honesta quanto aos pontos fortes e fracos. Em \textit{Security} e \textit{Maintainability} (notas de 4,5) o sistema atinge um nível bom: BCrypt com \textit{cost factor}~11, \textit{login} \textit{constant-time}, \textit{lockout} após 5 falhas, rotação de \textit{refresh tokens}, \textit{audit log} universal, separação rigorosa em camadas e \textit{suite} de testes que corre em três segundos. Em \textit{Performance Efficiency} e \textit{Usability} (notas de 3), o produto carece de evidência empírica: não foram efetuados testes de carga nem auditoria de acessibilidade, e o SRS não fixa SLAs verificáveis em CI. Nenhuma das lacunas identificadas compromete a operação prevista para a cadeia de cinco lojas do caso de estudo, mas marcam claramente os pontos onde o sistema precisaria de evoluir antes de uma colocação em produção comercial.

%--------------------------------------------------------------------------
\section{Trabalho Futuro}
%--------------------------------------------------------------------------

    A reflexão crítica feita na secção anterior permite identificar várias linhas de trabalho futuro, organizadas em três eixos: completar funcionalmente o sistema, reforçar a evidência de qualidade não funcional e amadurecer a operação e o ciclo de vida do \textit{software}.

    No \textbf{primeiro eixo, completude funcional}, deverão ser implementados os quatro casos de uso colocados fora do âmbito por decisão pragmática: UC07 (inventário manual), UC08 (sugestões automáticas de reposição), UC14 (exportação de relatórios em PDF/CSV) e UC18 (interface de consulta de \textit{logs} de auditoria). Os respetivos diagramas de sequência e DTOs já estão preparados no Capítulo~4, pelo que o esforço incremental é sobretudo de \textit{front-end} e de testes de aceitação adicionais.

    No \textbf{segundo eixo, evidência de qualidade não funcional}, três trabalhos devem ser conduzidos. Primeiro, introduzir testes de carga sobre os \textit{endpoints} críticos do POS (em particular \texttt{POST /api/v1/vendas}) usando ferramentas como \texttt{k6} ou \texttt{Bombardier}, e fixar SLAs quantitativos no SRS (por exemplo, p99 inferior a 200~ms para registar uma venda) que possam ser verificados em CI. Segundo, fazer uma auditoria formal de acessibilidade ao \textit{front-end} Blazor, com integração de \texttt{axe-core} nos testes automáticos, garantindo conformidade com WCAG~2.1 AA. Terceiro, fechar as lacunas pontuais de cobertura de código identificadas no Capítulo~6 (\textit{branches} alternativos do \texttt{ProcurementService} e testes para o \texttt{ReportingService}).

    No \textbf{terceiro eixo, maturidade operacional}, propõe-se a migração dos testes de integração de \texttt{EF~Core~InMemory} para \texttt{Testcontainers.MsSql}, eliminando a divergência entre os comportamentos do \textit{provider} em memória e do SQL~Server real (nomeadamente em \textit{constraints} de unicidade e \textit{cascading deletes}). A par disso, deve ser configurado um \textit{pipeline} de CI/CD (GitHub Actions ou equivalente) que execute toda a \textit{suite} de testes a cada \textit{commit}, gere o relatório de cobertura e imponha um limite mínimo (por exemplo, 70\,\%) como condição de \textit{merge}, em linha com as práticas DevOps descritas por Leite et al. (2019). A observabilidade em produção (métricas Prometheus, \textit{dashboards} Grafana, \textit{logs} estruturados Serilog enviados para um agregador) é outro passo necessário antes de uma operação real.

    Como evolução posterior, e caso o sistema venha a ser colocado em produção comercial, será necessário endereçar a conformidade fiscal completa que ficou deliberadamente fora do âmbito do protótipo académico (ADR-009): ATCUD, SAF-T, encadeamento por \textit{hash} e certificação pela Autoridade Tributária. Estas características têm implicações arquiteturais não triviais (numeração de faturas com integridade criptográfica, retenção de dados por dez anos, exportação periódica em formato SAF-T) e ficariam mais bem tratadas como uma etapa autónoma com revisão jurídica especializada.

    Por fim, sobre a metodologia em si, a equipa considera que a fórmula de desenvolvimento assistido por LLM com revisão humana final, tal como foi praticada aqui, pode ser sistematizada como um padrão reutilizável em projetos futuros. As lições mais importantes a transportar são a necessidade de fornecer ao LLM o contexto técnico real antes de pedir produção (interfaces, casos de uso, diagramas), a importância de descrever \textit{bugs} de forma cirúrgica para obter correções focadas, e o cuidado de não delegar no LLM o critério de aceitação de regras de negócio, que tem de ficar do lado humano. Estes princípios alinham-se com as recomendações de Chen et al. (2021) sobre Engenharia de \textit{Software} para sistemas assistidos por IA.
